\begin{proof}
\begin{enumerate}[label=\arabic*.,leftmargin=*]
\item \cite[Lemma~9.10(ii), p.~70, and Proposition~9.11,
pp.~70--71]{AC2026} give the evaluation divisor \(2S_2\) and
\(f_*\omega_{X/B}\simeq\mathcal O_B(p_0)\) for these same
marked fillings. Retain the explicit base identifications
\[
 h=(t-1)^{-1},\quad \operatorname{div}(h)=p_0-p_2,\qquad
 \eta_B=dt/(t-1)^2,\quad\operatorname{div}(\eta_B)=-2p_2.
\]
Tensoring evaluation with \(f^*\eta_B\), and multiplying the
base line by \(f^*h\), gives (CR1), since \(f^*p_2=4S_2\).

\item By \cref{inv:A1.0}, a section of
\(\mathcal O_X(kS_2)\), viewed as a meromorphic function,
descends away from \(p_2\) to a function \(b(t)\).
On the actual finite cover \(t-1=s_2^4\), its pullback is
invariant under \(s_2\mapsto\zeta_2s_2\), so its Laurent
exponents are multiples of four. The exact order condition is
\[
 4\operatorname{ord}_{p_2}b+k\ge0
 \quad\Longleftrightarrow\quad
 \operatorname{ord}_{p_2}b\ge-\lfloor k/4\rfloor.
\]
Conversely every such \(b\) pulls back to an allowed section.
This is (CR2) as an equality of subsheaves, and applying it
to \(k=2m,-2m\) gives (CR3).

\item Fix the isomorphism
\(\omega_X^{-1}\simeq\mathcal O_X(2S_2)\) from the first step.
Let \(u\) correspond to its canonical divisor section and
choose \(\ell\in H^0(B,\mathcal O_B(p_2))\) with
\(\ell(p_2)\ne0\). Let \(v\) correspond to \(f^*\ell\), so
\(u^2=f^*\sigma_{p_2}\) under the induced degree-two map.
The map \(U\mapsto u,V\mapsto v\) sends the degree-\(m\)
basis \(U^{m-2j}V^j\), \(0\le j\le\lfloor m/2\rfloor\),
to the pullback basis from (CR2), multiplied by \(u\) when
\(m\) is odd. It is therefore the graded isomorphism (CR4).
Replacing \(\ell\) by \(a\ell+b\sigma_{p_2}\) gives
\(v\mapsto av+bu^2\); multiplying the chosen line isomorphism
by \(c\ne0\) gives \((u,v)\mapsto(c^{-1}u,c^{-2}v)\).
The base-point-free pair \((u^2,v)\) induces exactly
\[
 x\longmapsto[\sigma_{p_2}(f(x)):\ell(f(x))].
\]
For \(\ell=A/(t-1)+B\), \(A\ne0\), the specified base map is
\(t\mapsto[t-1:A+B(t-1)]\). Degree one has the single
divisor \(2S_2\).
\end{enumerate}
\end{proof}
